% Cuerpo científico recuperado y distribuido en su residencia terminal.
% Cuerpo editor-ready sin preámbulo, portada ni metadatos publicables.
% Complementa, sin sustituirlos, los cierres operatoriales y espectrales.

\section{Clausuras negativas y calibradas de las equivalencias aparentes}
\label{sec:hmt-final-clausuras-negativas-calibradas}

Los resultados de esta sección se aplican después de la generación conjunta
del continuo. Su antecedente causal es el estado único producido por
APP--TRIT--TPK, la emergencia correlativa de sus cinco construcciones
consustanciales, los cuadrados de naturalidad y el paso conjunto al terminal
y al límite. Ninguna de las equivalencias examinadas a continuación genera o
selecciona retrospectivamente esas construcciones.

\subsection{Obstrucciones de canonicidad en la frontera
\texorpdfstring{\(108{:}144=90{:}120\)}{108:144=90:120}}

\begin{theorem}[No existe una inclusión canónica
\texorpdfstring{\(6\hookrightarrow8\)}{6 a 8} sin calibre]
\label{thm:hmt-final-no-inclusion-canonica-seis-ocho}
Sean \(A\) y \(B\) conjuntos desnudos con

\[
 |A|=6,
 \qquad |B|=8.
\]
No existe una regla, natural respecto de todas las relabelizaciones de \(A\)
y \(B\), que seleccione una inyección \(A\hookrightarrow B\).
\end{theorem}

\begin{proof}
Si una inyección \(\iota:A\hookrightarrow B\) fuese seleccionada sin dato
adicional, la naturalidad frente a todo \(\beta\in\operatorname{Sym}(B)\)
obligaría a \(\beta\iota=\iota\). Elegido \(a\in A\), existe una permutación
\(\beta\) que mueve \(\iota(a)\), contradicción. Por tanto, el transductor
ruta a ruta APP--Witt no es canónico a partir de los cardinales. Una vez
fijados una octada, una inclusión y el calibre de incidencia, sí se obtiene un
transductor calibrado; la elección forma parte de su dominio.
\end{proof}

\begin{theorem}[Una fibra desnuda de cardinal \(28\) no determina
\(J(8,2)\)]
\label{thm:hmt-final-no-johnson-desde-cardinal}
Sea \(F\) un conjunto sin estructura con \(|F|=28\). No existe una relación
de adyacencia no trivial sobre \(F\), natural respecto de todas sus
permutaciones, que produzca el grafo de Johnson \(J(8,2)\).
\end{theorem}

\begin{proof}
El grupo \(\operatorname{Sym}(F)\) es transitivo sobre los pares no ordenados
de elementos distintos. Una relación simple e invariante bajo todo ese grupo
contiene, por tanto, todos los pares o ninguno. Los dos grafos resultantes son
el completo y el vacío. En cambio, \(J(8,2)\) tiene \(28\) vértices y
grado \(12\); no es ninguno de ellos. Por consiguiente, la proyección
calibrada \(1008\to36\) con fibras uniformes de cardinal \(28\) no implica
por sí sola una octada de Witt. Fijados una octada \(O\) y una biyección
\(F\simeq\binom{O}{2}\), la adyacencia de Johnson puede transportarse a \(F\),
pero depende precisamente de ese calibre.
\end{proof}

Estos dos teoremas cierran negativamente las antiguas formulaciones
canónicas. No afectan al invariante exacto
\[
 \frac{108}{144}=\frac{90}{120}=\frac68
\]
ni al defecto normalizado \(-1/12\), que pertenecen a la extensión de rango
y no a la elección de una incidencia de Witt.

\subsection{La vuelta nonádica completa y el lector excepcional}

Fijemos el calibre
\[
 \chi_{\mathrm{exc}}
 =(
 \prec_{\mathrm{proj}},
 \kappa_{\mathrm{Paley}},
 \mathcal D_{12},
 \iota_{\mathrm{inc}},
 \mathcal K_{\Lambda})
\]
y el lector calibrado
\(E_{\chi_{\mathrm{exc}}}:\mathscr S_\infty
\to\mathcal I^{\mathrm{exc}}_{\chi_{\mathrm{exc}}}\).
Sea
\[
 \mathscr S_\infty\xleftarrow{s_\infty}
 \mathsf Z^{\Gamma}_{9,\infty}
 \xrightarrow{t_\infty}\mathscr S_\infty
\]
la arista testificada de una vuelta completa. Denotemos por
\(\operatorname{ed}_0\) la arista central, por \(\operatorname{ev}_0\) la
evaluación en el estado inicial y por \(F_9\) el desplazamiento de una vuelta.

\begin{theorem}[Factorización excepcional por la arista testificada]
\label{thm:hmt-final-factorizacion-gamma-nueve}
Defínanse
\[
 \overline E^{s}_{\chi_{\mathrm{exc}}}
 :=E_{\chi_{\mathrm{exc}}}\circ s_\infty,
 \qquad
 \overline E^{t}_{\chi_{\mathrm{exc}}}
 :=E_{\chi_{\mathrm{exc}}}\circ t_\infty.
\]
Entonces
\begin{align}
 E_{\chi_{\mathrm{exc}}}\circ\operatorname{ev}_0
 &=\overline E^{s}_{\chi_{\mathrm{exc}}}
   \circ\operatorname{ed}_0,
 \label{eq:hmt-final-factorizacion-gamma-fuente}\\
 E_{\chi_{\mathrm{exc}}}\circ\operatorname{ev}_0\circ F_9
 &=\overline E^{t}_{\chi_{\mathrm{exc}}}
   \circ\operatorname{ed}_0.
 \label{eq:hmt-final-factorizacion-gamma-destino}
\end{align}
\end{theorem}

\begin{proof}
Las identidades de extremos son
\(s_\infty\operatorname{ed}_0=\operatorname{ev}_0\) y
\(t_\infty\operatorname{ed}_0=\operatorname{ev}_0F_9\). La composición con
el lector calibrado da las dos ecuaciones.
\end{proof}

\begin{corollary}[No equivalencia entre el sector temporal \(+1\) y \(1\)
vuelta nonádica completa]
\label{cor:hmt-final-no-equivalencia-retorno-gamma}
La factorización anterior no puede reemplazarse, de forma natural, por una
factorización a través de la sola muestra TRIT \(\tau=+1\).
\end{corollary}

\begin{proof}
El valor \(\tau=+1\) registra retorno, memoria y pasado, pero olvida los datos
de fase, acarreo, ruta y extremo que distinguen aristas completas. Dos aristas
con la misma muestra temporal pueden tener destinos diferentes. Todo mapa que
factorice por esa muestra identifica ambas aristas, mientras
\(\operatorname{ed}_0\) y sus mapas fuente y destino las separan. Por tanto,
el sector \(+1\) no es equivalente a la arista completa \(K\mapsto K+9\).
En particular,
\[
 g(K+9)=g(K)
 \quad\text{y}\quad
 B_{K+9}\ne B_K\quad\text{en general}
\]
siguen siendo afirmaciones distintas.
\end{proof}

\begin{proposition}[Separación calibrada frente a publicación cruda]
\label{prop:hmt-final-separacion-enriquecida-no-cruda}
Supóngase que el lector excepcional enriquecido conserva, para cada
generador, el marco \(f_j\) y la imagen
\(y_j=\Gamma_{f_jA_Wf_j^{-1}}(w_j)\). Entonces recupera
\[
 w_j=\Gamma_{f_jA_Wf_j^{-1}}^{-1}(y_j)
\]
y es inyectivo en el cociente por los cambios de coordenadas incluidos en el
calibre. Si el lector crudo se obtiene olvidando marcos o memoria,
\(E^{\mathrm{cr}}=\pi E^{\sharp}\), no es equivalente al lector enriquecido
salvo que \(\pi\) sea inyectiva sobre \(E^{\sharp}(X)\).
\end{proposition}

\begin{proof}
La primera afirmación es la inversión explícita de un automorfismo en cada
fibra. Para la segunda, si \(\pi\) identifica dos salidas enriquecidas
distintas, el lector crudo identifica sus antecedentes y no puede ser
inyectivo. Recíprocamente, la inyectividad de \(\pi\) sobre la imagen permite
recuperar la salida enriquecida. Así se cierra la separación en el dominio
calibrado y se rechaza la equivalencia silenciosa con el codominio crudo.
\end{proof}

El antiguo cuadrado escrito con símbolos no tipados
\(\mathrm{pub}_{\mathrm{inc}},\mathrm{coord},\mathcal O\) no constituye una
equivalencia matemática: no define siquiera sus morfismos. Su sustituto es el
cono natural de las publicaciones arquimediana, excepcional calibrada y
dimensional, cuyos cuadrados de truncación sí tienen dominio, codominio y
acción declarados.

\subsection{Corredor excepcional y lazos del transporte}

La cadena Paley--Witt--Golay--Leech--\(V^{\natural}\)--Monstruo--Moonshine
es una publicación posterior del estado ya construido. Una estrella de
incidencia y un lazo del grupoide TPK no son objetos equivalentes por su
nombre: la primera es un subobjeto combinatorio; el segundo es una flecha con
fuente igual a destino, composición e inversa. Sin base, orientación y funtor
de transporte no existe una comparación tipada. Por ello, la identificación
automática queda cerrada como no-equivalencia. Un funtor calibrado de
transporte puede producir lazos a partir de estrellas, pero constituye un
resultado adicional y no una premisa del corredor ya demostrado.

\subsection{Canal infinitamente divisible inducido por la transferencia TPK}

\begin{theorem}[Poissonización UCP y raíces coherentes]
\label{thm:hmt-final-poissonizacion-ucp}
Sea \(K\) un canal UCP compatible con los truncamientos TPK\@. Para \(t\ge0\),
defínase
\[
 P_t:=\exp\bigl(t(K-I)\bigr)
 =e^{-t}\sum_{n=0}^{\infty}\frac{t^n}{n!}K^n.
\]
Entonces \((P_t)_{t\ge0}\) es un semigrupo UCP compatible con todos los
truncamientos y
\[
 P_1=(P_{1/m})^m
 \qquad(m\ge1).
\]
En cada nivel finito, una dilatación mínima de Stinespring de \(P_{1/m}\)
produce una raíz tensorial coherente; el límite proyectivo conserva las
marginales TPK.
\end{theorem}

\begin{proof}
Cada \(P_t\) es una combinación convexa normoconvergente de los canales
\(K^n\), por lo que es UCP\@. Como todos los términos son potencias del mismo
operador, la identidad exponencial da \(P_tP_s=P_{t+s}\). La compatibilidad
con los truncamientos se verifica término a término. Tomando \(t=1/m\) se
obtiene la raíz \(m\)-ésima. Stinespring convierte cada composición finita en
una contracción tensorial; la naturalidad de los canales hace coherentes sus
marginales.
\end{proof}

El teorema cierra la existencia de un canal infinitamente divisible con
memoria TPK\@. No afirma que un canal histórico arbitrario coincida con esta
Poissonización: esa igualdad requeriría comparar sus generadores y queda
rechazada mientras no se proporcione tal comparación.

\section{Clausuras espectrales y fronteras de reconocimiento externo}
\label{sec:hmt-final-clausuras-espectrales-externas}

\begin{proposition}[Equivalencia calibrada, no igualdad ambiental]
\label{prop:hmt-final-equivalencia-calibrada-vnueve}
El transporte espectral del doble círculo determina una clase unitaria en el
subespacio cíclico generado por su vector de referencia. Dos representantes
ambientales \(V_9\) son literalmente iguales sólo después de fijar una
identificación unitaria del complemento y un calibre compatible. Sin esos
datos, la igualdad ambiental no queda indeterminada: es una afirmación más
fuerte y no equivalente a la clase específica del subespacio ya construida.
\end{proposition}

La convergencia en norma de los resolventes comprimidos, la recuperación de
alturas, el criterio de Riesz para multiplicidades y la traza diagonal
RH--Moonshine quedan cerrados en sus teoremas respectivos. La traza diagonal
acopla el entero de memoria con el grado de \(V^{\natural}\); no afirma que la
multiplicidad de cada cero sea la dimensión de una representación del
Monstruo. Esa identificación adicional no es equivalente a la existencia de
la traza acoplada.

Quedan fuera del dominio interno, sin reabrir ningún teorema HMT:
\begin{enumerate}
 \item la identificación de la dicotomía de publicaciones HMT analíticas con
 todos los subconjuntos arbitrarios de \(\mathbb R\) del universo ambiente;
 \item la realización macroscópica o termodinámica del cristal temporal
 finito HMT ya construido, incluida la conservación de coherencia a gran
 volumen y su contraste experimental;
 \item el endurecimiento decimal de ventanas espectrales mediante una
 ejecución concreta de aritmética intervalar;
 \item una identificación representacional adicional entre multiplicidades
 espectrales individuales y módulos irreducibles del Monstruo.
\end{enumerate}
Estas cuatro fronteras son reconocimientos o comparaciones externos, no
lagunas de la construcción del continuo, del cierre cardinal, de RH ni del
entrelazamiento del doble círculo.

La residencia terminal conserva, sin permutación,
\[
\begin{aligned}
 \mathrm{RH}
 &\longrightarrow \text{doble círculo y campo espectral}
 \longrightarrow \mathrm{Navier\text{--}Stokes}\\
 &\longrightarrow \mathrm{Yang\text{--}Mills}
 \longrightarrow \mathrm{Hodge}
 \longrightarrow \mathrm{BSD}
 \longrightarrow \text{falsadores}.
\end{aligned}
\]
